\section{Analyse der bestehenden Szenarien im Mehrsinnlichen Interaktionsraum (MIR)}
\label{sec:analyse}
Im Rahmen einer explorativen Analyse wurden bestehende Anwendungsszenarien des projektionsbasierten Multisensorischen Interaktionsraums (MIR) an der HTW Dresden systematisch erfasst und hinsichtlich zentraler Gestaltungsdimensionen untersucht.

Ziel dieser Analyse war es, strukturelle Muster der Szenariengestaltung, technologische Konfigurationen sowie Interaktionsparadigmen zu identifizieren. Auf dieser Grundlage sollten Anforderungen für die Entwicklung körper\-basierter Interaktionskonzepte im MIR abgeleitet werden.

Die zugrunde liegenden Analyse- und Kategorisierungsdimensionen wurden induktiv, also anhand von Beobachtungen und Daten der bestehenden Szenarien, entwickelt. Sie umfassen unter anderem Inhaltstypen, Navigationskonzepte, physische Inszenierungselemente sowie vorhandene Interaktionsformen.
% induktiv = ich schaue mir szenarien an und erkenne dann das Muster statt ich nehme Kategorien aus der Literatur und prüfe sie (das wäre deduktiv)

Die Analyse basiert auf einer Auswahl der aktuell verfügbaren MIR-Szenarien, die zum Zeitpunkt der Untersuchung im System implementiert und zugänglich waren. Insgesamt wurden neun Szenarien betrachtet, die unterschiedliche Anwendungskontexte wie atmosphärische Erlebnisräume, narrative Präsentationen sowie simulations- oder lehrorientierte Anwendungen abdecken.
Es ist zu berücksichtigen, dass im MIR darüber hinaus weitere Szenarien existieren können, die jedoch im Rahmen dieser Arbeit nicht systematisch untersucht wurden. Die betrachteten Szenarien stellen somit eine exemplarische Auswahl der aktuell genutzten Anwendungen dar und dienen als Grundlage für die anschließende Analyse der Interaktionsstrukturen.

Die untersuchten MIR-Szenarien befassen sich mit den folgenden Inhalten:
\begin{itemize}

    \item \textbf{Waldbaden:} Atmosphärisches Szenario, das vorab aufgezeichnete Video- und Audioaufnahmen einer natürlichen Waldumgebung nutzt.
    \item \textbf{Fußballstadion:} Atmosphärisches Szenario, das vorab aufgezeichnete Video- und Audioaufnahmen eines Fußballstadions mit mehreren perspektivischen Ansichten des Spielgeschehens nutzt.
    \item \textbf{Kokon:} Atmosphärisches Szenario, das vorab aufgezeichnete Video- und Audioaufnahmen als Innenansicht eines Kokons als abstrahierte Umgebung inszeniert.
    \item \textbf{Zug fahren:} Atmosphärisches Szenario, das vorab aufgezeichnete Video- und Audioaufnahmen mit wechselnden Ansichten aus einer Zugfahrt und situativen Haltepunkten zur Betrachtung der Umgebung nutzt.
    \item \textbf{YouTube 360°:} Wiedergabe externer 360°-Videoinhalte, die auf die Projektionsflächen des MIR angepasst werden.
    \item \textbf{Dschungel Level:} Mit der 3D-Software Unity live gerendertes atmosphärisches Szenario einer virtuellen Dschungelumgebung.
    \item \textbf{Smart Factory:} Live gerenderte Simulation einer Smart Factory mit digital repräsen\-tier\-ten Maschinen- und Prozessabläufen.
    \item \textbf{OP:} Lehrszenario, das vorab aufgezeichnete Video- und Audioaufnahmen einer realen Operationsumgebung zu Simulations- und Demonstrationszwecken nutzt.
    \item \textbf{Hubschrauberlandeplatz:} Lehrszenario, das vorab aufgezeichnete Video- und Audioaufnahmen eines realen Hubschrauberlandeplatzes zu Simulations- und Demonstrationszwecken nutzt.
    \item \textbf{(KI-Szenarien):} Konzeptuell vorgesehene, KI-generierte Umgebungen, die durch Texteingaben der Nutzer:innen verändert werden können.
\end{itemize}


\subsection{Inhaltstypen und mediale Ausgestaltung}
Die untersuchten Szenarien lassen sich hinsichtlich ihrer medialen Inhaltsstruktur in drei übergeordnete Typen einteilen:

\begin{enumerate}
\item vorgerenderter Inhalt
\item live gerenderter Inhalt
\item KI-generierter Inhalt
\end{enumerate}

Der Großteil der bestehenden Anwendungen basiert auf vorgerenderten Inhalten.
Hierbei handelt es sich in der Regel um vorab aufgezeichnete Video- und Audioaufnahmen realer Situationen, die für die Projektion auf die 270°-Leinwände räumlich angepasst und synchronisiert wurden \parencite{orawetz2024mir}.
Live gerenderte Inhalte werden in der Unreal Engine generiert und ermög\-lichen eine direkte \textit{Manipulation} oder Navigation innerhalb der virtuellen Szene.
KI-generierte Inhalte wurden bislang nur punktuell berücksichtigt und weisen noch keinen integrativen Stellenwert innerhalb der bestehenden Szenarien auf. Diese Inhalte zielen darauf ab, Veränderungen der Umgebung über Texteingaben der Nutzer:innen zu ermög\-lichen.

Technologisch basieren die meisten Anwendungen auf einer Kombination aus 270°-Visuali\-sierung, 360°-Audiosystem und teilweise zusätzlichem Bodenfeedback (Vibration) (siehe Tabelle \ref{tab:charakteristika}). Diese multisensorische Konfiguration dient primär der Immersionssteigerung des Raumerlebens. Die technische Infrastruktur erscheint damit vor allem auf Wahrnehmungsintensivierung und atmosphärische Verdichtung ausgerichtet, während interaktive Dynamiken in den bestehenden Szenarien bislang weniger stark ausgeprägt sind.


\begin{table}[H]
\centering
\small
\caption{Verdichtete Charakteristika der analysierten MIR-Szenarien}
\begin{tabular}{p{4cm} p{3cm} p{5cm} p{3cm}}
\toprule
\textbf{Szenario} & \textbf{Inhaltstyp} & \textbf{Verwendete Technologie} & \textbf{Szenarioziel} \\
\midrule
Waldbaden & vorgerendert & 360° Sound, 270° Visuals & atmosphärisch (ambient) \\
Fußballstadion & vorgerendert & 360° Sound, 270° Visuals, Bodenvibration & atmosphärisch (narrativ) \\
Kokon & vorgerendert & 360° Sound, 270° Visuals, Bodenvibration & atmosphärisch (ambient) \\
Zug fahren & vorgerendert & 270° Visuals & atmosphärisch (ambient) \\
YouTube 360° & vorgerendert & 360° Sound, 270° Visuals & atmosphärisch (narrativ) \\
Jungle Level & live gerendert & 360° Sound, 270° Visuals & atmosphärisch (ambient) \\
Smart Factory & live gerendert & 360° Sound, 270° Visuals & Simulation \\
OP & vorgerendert & 360° Sound, 270° Visuals & Lehre \\
Hubschrauberlandeplatz & vorgerendert & 360° Sound, 270° Visuals & Lehre \\
\bottomrule
\end{tabular}
\label{tab:charakteristika}
\end{table}

Eine detaillierte Übersicht aller erfassten Szenarien und Analysekategorien befindet sich im Anhang (siehe Abbildung~\ref{tab:scenario}).


\subsection{Physische Raumgestaltung}
Ein weiteres Merkmal einiger Szenarien ist der Einsatz physischer Inszenierungselemente zur atmosphärischen Anreicherung. Hierzu zählen unter anderem Wiesenboden, Stühle, textile Elemente, ein manueller Reparaturarbeitsplatz  oder Lehrobjekte. Die meisten dieser materiellen Ergänzungen fungieren primär als haptisch-visuelle Verstärker der dargestellten Szenarien und unterstützen die Kohärenz zwischen physischem und projiziertem Raum. Eine Ausnahme bildet das Szenario \textit{Smart Factory}, in dem der manuelle Reparaturarbeitsplatz funktional mit dem digitalen System gekoppelt ist und dadurch nicht nur atmosphärisch, sondern auch interaktiv wirksam wird.
Es zeigt sich jedoch, dass die meisten dieser Objekte primär dekorativ oder symbolisch eingesetzt werden. Eine systematische sensorische Einbindung (zum Beispiel über NFC, UWB oder andere Tracking-Technologien) zur Erzeugung interaktiver Zustandsänderungen ist bislang kaum implementiert. Somit bleibt das Potenzial physischer Objekte als Schnittstelle zwischen Nutzer:in und digitalem System weitgehend ungenutzt.


\subsection{Szenarioziele}
Die Mehrheit der Szenarien verfolgt primär ein atmosphärisches Ziel. Im Vordergrund steht dabei nicht die aktive Bearbeitung eines Systems, sondern das Eintauchen in eine medial inszenierte räumliche Situation. Innerhalb der atmosphärischen Szenarien lassen sich ambient orientierte und narrativ strukturierte Formen unterscheiden. Zu den ambienten Szenarien zählen beispielsweise der Kokon und der Wald, da dort die Nutzer:innen nicht an bestimmte zeitliche Ereignisse gebunden sind und die Umgebung eher auf ein kontinuierliches atmosphärisches Erleben ausgerichtet sind. Ambiente atmosphärische Szenarien verfolgen also primär das Ziel, eine immersive Stimmung zu erzeugen, ohne dass eine aktive Interaktion erforderlich ist. In solchen Fällen könnte zusätzliche Interaktivität sogar als störend wahrgenommen werden. Narrative Szenarien hingegen, wie beispielsweise das Fußballstadion, folgen einer zeitlich strukturierten Ereignisabfolge. Hier könnte Interaktion potenziell genutzt werden, um Übergänge, Perspektivwechsel oder Zusatzinformationen zu steuern.

Simulations- und Lehrszenarien weisen tendenziell höhere Interaktivitätsgrade auf, allerdings erfolgt die Steuerung bisher meist über klassische Eingabegeräte am Rechner oder über explizite User-Interface-Elemente (UI Elemente). Die Interaktion bleibt somit gerätezentriert.


\subsection{Navigationslogik und Interaktionsformen}
Die Navigation innerhalb der untersuchten Szenarien stellt eine grundlegende Form der Interaktion dar und erfolgt überwiegend systemgesteuert über fest definierte Perspektivwechsel in vorgerenderten Sequenzen. Diese Perspektivwechsel sind typischerweise videobasiert und an vorab definierte Sprungpunkte gebunden.

In einzelnen Fällen, insbesondere bei live-gerenderten Szenarien, wird eine freie Navigation ermöglicht. Diese erfolgt jedoch meist über klassische Eingabegeräte wie Tastatur und Maus. Die Navigation ist damit primär systemgetrieben und nur eingeschränkt nutzerzentriert. Zwar können sich Nutzer:innen physisch frei im Raum bewegen, diese Bewegung ist jedoch in der Regel nicht mit einer Veränderung des medialen Zustands gekoppelt. Eine direkte Verbindung zwischen Körperbewegung im Raum und inhaltlicher Systemreaktion ist somit bislang kaum realisiert.

Über die Navigation hinaus zeigt sich insgesamt eine geringe Ausprägung systemischer Reaktivität. In vielen Szenarien ist keine Interaktion vorgesehen; in anderen erfolgt sie ausschließlich explizit über Rechner, Controller oder UI-Schaltflächen.

Wo Interaktion implementiert ist, geschieht sie überwiegend in folgenden Formen:
\begin{itemize}
\item UI-basierte Eingaben (zum Beispiel Schaltflächen oder Prompteingabe)
\item Objektinteraktionen (zum Beispiel im Szenario Smart Factory (siehe Abschnitt \ref{sec:rw_multisensory}))
\item Tastatur- und Maussteuerung zur Navigation in virtuellen Umgebungen
\end{itemize}

Diese Interaktionen sind meist diskret, explizit und funktional motiviert. Eine kontinuierliche, implizite Kopplung zwischen Körperzustand und Systemverhalten ist kaum ausgeprägt. Darüber hinaus reagiert das System in den meisten Fällen nicht auf alltägliche oder subtile Bewegungen. Die Interaktionslogik ist somit überwiegend ereignisbasiert statt kontinuierlich zustandsbasiert.

Zur systematischen Gegenüberstellung der Szenarien werden die beobachteten Interaktionsstrukturen entlang mehrerer Dimensionen erfasst. Die Interaktionsform beschreibt dabei die Art der Eingabe, über die Nutzer:innen mit dem System interagieren (zum Beispiel keine Interaktion, UI-basierte Eingaben oder Tastatur-/Maussteuerung).
Die Systemreaktion bezeichnet die durch diese Eingaben ausgelöste Veränderung des Systemzustands, etwa Perspektivwechsel oder Zustandsänderungen innerhalb der Szene.
Der Interaktionstyp unterscheidet, ob Interaktionen explizit oder implizit erfolgen.
Die Navigation beschreibt, wie sich Nutzer:innen innerhalb des Szenarios bewegen beziehungsweise wie Perspektivwechsel realisiert werden.

Die unterschiedlichen Ausprägungen dieser Dimensionen sind in Tabelle~\ref{tab:interaktion} zusammengefasst.

\begin{table}[H]
\centering
\small
\caption{Interaktionsformen und Systemreaktivität im MIR}
\begin{tabular}{p{2,5cm} p{3cm} >{\raggedright\arraybackslash}p{3cm} p{1,5cm} >{\raggedright\arraybackslash}p{4cm}}
\toprule
\textbf{Szenario} & \textbf{Interaktions\-form} & \textbf{System\-reaktion} & \textbf{Inter\-aktions\-typ} & \textbf{Navigation}\\
\midrule
Waldbaden & keine & keine & keine & fester Perspektiv\-wechsel \\
Fußballstadion & keine & keine & keine & fester Perspektiv\-wechsel \\
Kokon & keine & keine & keine & keine \\
Zug fahren & keine & zeitbasierter Wechsel & keine & fester Perspektiv\-wechsel \\
YouTube 360° & keine & keine & keine & fester Perspektiv\-wechsel \\
Jungle Level & Tasta\-tur/Maus & Bewegung im Raum & explizit & frei (mit Tasta\-tur) \\
Smart Factory & Tasta\-tur/Maus, Gesteninteraktion & Zustands\-än\-der\-ung im Re\-pa\-ra\-tur\-ar\-beits\-platz ändert System\-zu\-stän\-de & explizit & eingeschränkt (mit Tasta\-tur) \\
OP & keine & keine & keine & fester Perspektiv\-wechsel \\
Hubschrauber\-landeplatz & keine & keine & keine  & keine \\
\bottomrule
\end{tabular}
\label{tab:interaktion}
\end{table}

Insgesamt zeigt sich, dass sowohl Navigation als auch Interaktion bisher überwiegend system- oder gerätebasiert umgesetzt wurden und kaum an die körperliche Präsenz der Nutzer:innen im Raum gekoppelt sind.


\subsection{Räumlich-körperliche Einflussfaktoren}
Darüber hinaus wurde im Rahmen der Untersuchung auch die räumliche Nutzung des Systems beobachtet. Dabei zeigte sich, dass körperliche Faktoren einen wesentlichen Einfluss auf das Immersionserleben haben. 
Beispielsweise verändert sich die visuelle Wahrnehmung der Projektionen abhängig von der Körperhaltung der Nutzer:innen. Bei Nutzung der Bodenvibrationsplatte im Stehen verschiebt sich der Blickwinkel tendenziell nach oben, wodurch Teile der Projektion außerhalb des direkten Sichtfeldes liegen können.
Darüber hinaus ist aufgrund der drei Projektionsflächen eine bevorzugte Blickrichtung im Raum gegeben. Die immersive Wahrnehmung ist daher nicht in allen Raumrichtungen gleichermaßen ausgeprägt.

Diese Beobachtungen verdeutlichen, dass körperliche Präsenz und räumliche Konfiguration untrennbar mit dem Interaktionsdesign verbunden sind. Gleichzeitig werden diese Faktoren bislang eher als atmosphärische Gegebenheiten statt als aktive Steuerungsdimension verstanden.


\subsection{Zusammenfassende Bewertung}
Die Analyse zeigt, dass bestehende MIR-Szenarien primär als immersive, multisensorische Erfahrungsräume konzipiert sind, jedoch wenige interaktive Möglichkeiten aufweisen.
Die bestehende Interaktionsarchitektur lässt sich überwiegend wie folgt charakterisieren:
\begin{itemize}
\item passiv-rezeptiv statt partizipativ, % Nutzer konsumiert nur Inhalte / schaut nur zu & beeinflusst nicht Ablauf (das wäre partizipativ)
\item explizit statt implizit,
\item gerätezentriert statt körperzentriert,
\item ereignisgesteuert statt kontinuierlich zustandsbasiert. % reagiert auf Events; kontinuierlich Zustandsbsaiert wäre System reagiert auf kontinuierliche Zustände zB positionsbasierte Interaktion
\end{itemize}

Gleichzeitig zeigt sich, dass die technische Infrastruktur des MIR (besonders räumliche Bewegungsfreiheit, multisensorische Ausgabesysteme und potenzielle Sensorintegration) grund\-sätzlich geeignete Voraussetzungen für körper\-basierte Interaktionsformen bietet.

Die zuvor beschriebenen Beobachtungen lassen sich zu übergeordneten Mustern verdichten (siehe Tabelle \ref{tab:muster}).

\begin{table}[H]
\centering
\small
\caption{Abgeleitete Muster der MIR-Szenarien}
\begin{tabular}{>{\raggedright\arraybackslash}p{5cm} >{\raggedright\arraybackslash}p{7cm}}
\toprule
\textbf{Dimension} & \textbf{Beobachtung} \\
\midrule
Inhalt & überwiegend vorgerendert \\
\midrule
Navigation & primär systemgesteuert, vordefiniert \\
\midrule
Interaktion & selten, meist explizit und gerätebasiert \\
\midrule
Systemreaktivität & kaum vorhanden, ereignisbasiert \\ %kaum vorhanden → Nutzer beeinflussen das System fast nicht; ereignisbasiert → nur bei expliziten Aktionen, nicht kontinuierlich
\midrule
Körperintegration & nicht mit Systemzustand gekoppelt \\
\midrule
Physische Inszenierungselemente & meist dekorativ, kaum interaktiv eingebunden \\
\bottomrule
\end{tabular}
\label{tab:muster}
\end{table}

% Übergang
Damit wird ein Spannungsfeld zwischen der stark immersiven räumlichen Gestaltung und der bislang nur begrenzt vorhandenen Interaktivität sichtbar.
Im folgenden Kapitel werden daher mögliche körper\-basierte Interaktionsformen untersucht und hinsichtlich ihrer Eignung für den Einsatz im MIR bewertet.